% !TEX program = xelatex
\documentclass[10pt, a4paper]{article}

\usepackage{fontspec}
\usepackage{polyglossia}
\setmainlanguage{russian}
\setotherlanguage{english}
\setmainfont{Times New Roman}
\setsansfont{Helvetica Neue}
\setmonofont{Menlo}
\newfontfamily\cyrillicfonttt{Times New Roman}

\usepackage[top=2cm,bottom=2cm,left=2.2cm,right=2.2cm]{geometry}
\usepackage{xcolor}
\usepackage{titlesec}
\usepackage{enumitem}
\usepackage{booktabs}
\usepackage{array}
\usepackage{longtable}
\usepackage{tabularx}
\usepackage{multirow}
\usepackage{fancyhdr}
\usepackage{amsmath}
\usepackage{hyperref}
\usepackage{tcolorbox}
\usepackage{colortbl}
\usepackage{parskip}
\usepackage{makecell}
\usepackage{pdflscape}

\tcbuselibrary{skins, breakable}

% ─── Colors ─────────────────────────────────────────────────────────────────
\definecolor{navy}{HTML}{1A237E}
\definecolor{steelblue}{HTML}{0277BD}
\definecolor{teal}{HTML}{00897B}
\definecolor{amber}{HTML}{E65100}
\definecolor{lightblue}{HTML}{E3F2FD}
\definecolor{lightgreen}{HTML}{E8F5E9}
\definecolor{lightyellow}{HTML}{FFF9C4}
\definecolor{lightorange}{HTML}{FFE0B2}
\definecolor{lightred}{HTML}{FFCDD2}
\definecolor{lightgem}{HTML}{E1F5EE}
\definecolor{lightstar}{HTML}{FFF8E1}
\definecolor{darkgray}{HTML}{333333}
\definecolor{midgray}{HTML}{607D8B}
\definecolor{rowalt}{HTML}{F5F5F5}

% ─── Section styling ────────────────────────────────────────────────────────
\titleformat{\section}{\large\bfseries\color{navy}}{\thesection.}{0.5em}{}
  [{\color{steelblue}\titlerule[1pt]}]
\titleformat{\subsection}{\normalsize\bfseries\color{steelblue}}{\thesubsection.}{0.4em}{}
\titleformat{\subsubsection}{\small\bfseries\color{midgray}}{\thesubsubsection.}{0.4em}{}

% ─── Boxes ──────────────────────────────────────────────────────────────────
\newtcolorbox{gembox}[1]{
  colback=lightgem, colframe=teal, coltitle=white,
  fonttitle=\bfseries\small, title={💎 #1},
  arc=4pt, left=6pt, right=6pt, top=5pt, bottom=5pt, breakable
}
\newtcolorbox{starbox}[1]{
  colback=lightstar, colframe=amber, coltitle=white,
  fonttitle=\bfseries\small, title={★ #1},
  arc=4pt, left=6pt, right=6pt, top=5pt, bottom=5pt, breakable
}
\newtcolorbox{stratbox}[1]{
  colback=lightblue, colframe=navy, coltitle=white,
  fonttitle=\bfseries, title={#1},
  arc=4pt, left=8pt, right=8pt, top=6pt, bottom=6pt, breakable
}
\newtcolorbox{tipbox}{
  colback=lightyellow, colframe=amber,
  arc=3pt, left=6pt, right=6pt, top=4pt, bottom=4pt, breakable
}

% ─── Header/Footer ──────────────────────────────────────────────────────────
\pagestyle{fancy}
\fancyhf{}
\fancyhead[L]{\small\color{darkgray}\textit{PhD Guide: Robotics / Autonomous Systems / Safe RL}}
\fancyhead[R]{\small\color{darkgray}\textit{Полный справочник профессоров}}
\fancyfoot[C]{\small\color{darkgray}\thepage}
\renewcommand{\headrulewidth}{0.4pt}

% ─── Helpers ────────────────────────────────────────────────────────────────
\newcommand{\match}[1]{%
  \ifnum#1>89 \colorbox{lightgreen}{\textcolor{teal}{\textbf{#1\%}}}%
  \else\ifnum#1>79 \colorbox{lightyellow}{\textcolor{amber}{\textbf{#1\%}}}%
  \else\ifnum#1>69 \colorbox{lightorange}{\textcolor{amber}{#1\%}}%
  \else \colorbox{lightred}{\textcolor{red!70!black}{#1\%}}%
  \fi\fi\fi}

\newcommand{\profsec}[3]{% {name}{lab}{match}
  \vspace{4pt}
  {\bfseries\color{navy}#1}\quad{\small\color{midgray}#2}\hfill\match{#3}\\[-2pt]
  {\color{steelblue}\rule{\linewidth}{0.3pt}}
}

\hypersetup{colorlinks=true,linkcolor=steelblue,urlcolor=teal}

% ════════════════════════════════════════════════════════════════════════════
\begin{document}
% ════════════════════════════════════════════════════════════════════════════

% ─── Title ──────────────────────────────────────────────────────────────────
\begin{center}
{\color{navy}\LARGE\bfseries PhD Guide: Robotics \& Autonomous Systems}\\[6pt]
{\color{steelblue}\large Полный справочник профессоров и стратегия поступления}\\[4pt]
{\color{navy}\rule{\linewidth}{2pt}}\\[4pt]
{\small\color{darkgray}
Темы интересов: Аэродинамика · Дроны (UAV) · Reinforcement Learning ·
System Control · Алгоритмы · Robotics \\
Ориентир: \textbf{Prof.\ Mykel Kochenderfer (Stanford SISL)} · 4-летний горизонт
}
\end{center}

\vspace{6pt}
\tableofcontents
\newpage

% ════════════════════════════════════════════════════════════════════════════
\section{Обозначения и легенда}
% ════════════════════════════════════════════════════════════════════════════

\begin{tabular}{lll}
\toprule
\textbf{Символ} & \textbf{Значение} & \textbf{Пояснение} \\
\midrule
\colorbox{lightgreen}{\textcolor{teal}{\textbf{90–100\%}}} & Исключительное совпадение &
  Темы лабы = твой список интересов \\
\colorbox{lightyellow}{\textcolor{amber}{\textbf{80–89\%}}} & Очень хорошее совпадение &
  Большинство тем совпадает \\
\colorbox{lightorange}{\textcolor{amber}{70–79\%}} & Хорошее совпадение &
  Частичное совпадение, стоит подавать \\
\colorbox{lightred}{\textcolor{red!70!black}{< 70\%}} & Умеренное совпадение &
  Запасной вариант \\
\colorbox{lightstar}{\textcolor{amber}{★}} & Рекомендую & Лучшие варианты в своём уровне \\
\colorbox{lightgem}{\textcolor{teal}{💎}} & Скрытая жемчужина &
  Сильный проф + менее конкурентный входной барьер \\
\bottomrule
\end{tabular}

% ════════════════════════════════════════════════════════════════════════════
\section{УРОВЕНЬ 1 — Топ-5 университеты}
% ════════════════════════════════════════════════════════════════════════════

\subsection{Stanford University}

% ─── Kochenderfer ────────────────────────────────────────────────────────────
\begin{starbox}{Mykel Kochenderfer — SISL (Stanford Intelligent Systems Lab) \hfill \match{98}}
\textbf{Темы:} MDP/POMDP, Safe RL, Collision avoidance для дронов и автопилотов, Julia/POMDPs.jl, Decision making under uncertainty, Верификация нейросетей, LLM safety.

\textbf{Почему подходит:}
ГЛАВНАЯ ЦЕЛЬ. Его темы — это буквально твой список интересов слово в слово. Автор трёх учебников MIT Press, четвёртый на подходе. Проект PID vs RL с POMDP-расширением — прямой разговор на его языке. Kochenderfer — единственный человек, где совпадение тем выходит за пределы 95\%.

\begin{tipbox}
\textbf{Совет:} Email писать после публикации первого эксперимента или arXiv-препринта. Обязательно упомянуть POMDPs.jl и Julia. Показать знакомство с его книгой «Algorithms for Decision Making».
\end{tipbox}
\end{starbox}

\vspace{4pt}

\profsec{Marco Pavone}{ASL — Autonomous Systems Lab}{82}
\textbf{Темы:} Autonomous aerospace vehicles, Control theory, Reachability, Safety, Motion planning, Optimization. Директор ASL, частично в NVIDIA.

\textbf{Почему:} Хорошая альтернатива если у Kochenderfer нет мест. Работает на autonomous aerospace vehicles и safety-critical control. Публикации с CBF и reachability.

\textbf{Совет:} Следить за его публикациями о reachability и motion planning. Писать про safety-critical control для autonomous systems.

% ─── MIT ─────────────────────────────────────────────────────────────────────
\subsection{MIT — Massachusetts Institute of Technology}

\begin{starbox}{Jonathan How — ACL (Aerospace Controls Lab) \hfill \match{90}}
\textbf{Темы:} UAV/дроны (физические), Multi-agent collision avoidance, POMDP-like planning, Neural networks + оптимизаторы, Robust planning, Distributed control.

\textbf{Почему:} Работает ИМЕННО на дронах и multi-vehicle collision avoidance. Нейросети + оптимизаторы на бортовых системах. Очень активная лаба, публикации в ICRA каждый год. Если не Stanford — MIT/How это второй приоритет.

\begin{tipbox}
\textbf{Совет:} Изучить его публикации 2024–2025 о safe autonomy и trajectory planning для UAV. В email указать конкретную статью.
\end{tipbox}
\end{starbox}

\vspace{4pt}

\profsec{Sertac Karaman}{LIDS — Lab for Information \& Decision Systems}{78}
\textbf{Темы:} UAVs, Stochastic systems, Алгоритмы (RRT*, sampling-based planning), Formal methods, Embedded systems. Директор LIDS.

\textbf{Почему:} Автор знаменитого алгоритма RRT*. Алгоритмическая база + дроны + случайные процессы. Хорошо если интересна математика алгоритмов planning.

\vspace{4pt}

\profsec{Navid Azizan}{LIDS / MechE}{72}
\textbf{Темы:} Meta-learning, Adaptive control, Дроны, Uncertainty, Mirror descent. Молодой профессор.

\textbf{Почему:} Недавняя статья о control под uncertainty именно для UAV с meta-learning. Молодой проф = больше шансов быть принятой, больше внимания студентам.

% ─── CMU ─────────────────────────────────────────────────────────────────────
\subsection{Carnegie Mellon University (CMU)}

\begin{starbox}{Guanya Shi — LeCAR (Learning and Control for Agile Robotics) \hfill \match{88}}
\textbf{Темы:} RL + control theory, Aerial robotics, Safe systems, Agile flight, Sim-to-real transfer.

\textbf{Почему:} PhD Caltech 2022, строит лабораторию с 2023 — много открытых позиций. Буквально: RL + control + aerial robotics + safety. Твой проект PID vs RL — прямое пересечение с его исследовательской программой.

\begin{tipbox}
\textbf{Совет:} Написать после 3-дневного прототипа. Показать код на GitHub. Молодой проф = больше личного внимания.
\end{tipbox}
\end{starbox}

\vspace{4pt}

\profsec{Sebastian Scherer}{AirLab — Aerial Robotics Lab}{75}
\textbf{Темы:} Autonomous aerial vehicles, Perception, Navigation, Field robotics, Sim-to-real.

\textbf{Почему:} Директор AirLab, фокус на физических дронах. Сильная практическая сторона. Подходит если интересен hardware + software вместе.

% ─── Caltech ─────────────────────────────────────────────────────────────────
\subsection{Caltech}

\profsec{Aaron Ames}{CAST — Center for AEroSpace Autonomy Research}{80}
\textbf{Темы:} Safety-critical systems, CBF (Control Barrier Functions), Nonlinear control, Robotics. Bren Professor.

\textbf{Почему:} Создатель CBF — математического фундамента safe RL. Если хочешь глубоко понять safety constraints, это самый правильный человек в мире. Очень теоретически сильная лаба.

\profsec{Soon-Jo Chung}{ARCL — Autonomous Robotics \& Control Lab}{77}
\textbf{Темы:} Aerospace robotics, Learning-based control, Multi-agent, Space/UAV, Estimation.

\textbf{Почему:} Научный руководитель Guanya Shi (CMU). Дроны + космос + learning control. Выпускники идут в JPL, Google, CMU.

% ─── Berkeley ────────────────────────────────────────────────────────────────
\subsection{UC Berkeley}

\profsec{Claire Tomlin}{HiPeR — Hybrid and Embedded Systems}{76}
\textbf{Темы:} Reachability analysis, Safe RL, Hybrid systems, UAV, Hamilton-Jacobi.

\textbf{Почему:} Пионер в reachability analysis для безопасного управления. Лучший выбор по formal safety в Berkeley. Изучить HJ-reachability перед контактом.

\profsec{Sergey Levine}{RAIL — Robotic AI \& Learning Lab}{65}
\textbf{Темы:} Deep RL, Robot learning, Imitation learning, Offline RL.

\textbf{Почему:} Топ в deep RL для роботов. Фокус на manipulation, меньше на UAV. Хорош как второй советник.

% ════════════════════════════════════════════════════════════════════════════
\section{УРОВЕНЬ 2 — Сильные программы}
% ════════════════════════════════════════════════════════════════════════════

\subsection{Georgia Tech}

\begin{gembox}{Shreyas Kousik — Safe Robotics Lab \hfill \match{85}}
\textbf{Темы:} Safe RL, Reachability analysis, Provably-safe neural networks, Aerospace applications.

\textbf{Почему:} Молодой (2022) профессор — активно берёт студентов. Специализация: safe RL через reachability + neural network verification. Статьи в ICRA 2025 о safe policy execution для роботов. Прямое совпадение с темой safety для дронов.

\textbf{Совет:} Изучи RAIL (Reachability-Aided Imitation Learning) — его последняя статья. Напиши про интерес к safe RL + UAV.
\end{gembox}

\profsec{Ye Zhao}{LIDAR Lab}{72}
\textbf{Темы:} RL + trajectory optimization, STL-MPC, Safe locomotion, Constrained control. Меньше фокус на дронах, больше на шагающих роботах, но методы те же.

% ─── Michigan ────────────────────────────────────────────────────────────────
\subsection{University of Michigan}

\begin{gembox}{Dimitra Panagou — DASC Lab (Distributed Autonomous Systems \& Control) \hfill \match{92}}
\textbf{Темы:} Safe multi-agent UAV системы, CBF, Provably-correct гарантии, Multi-vehicle coordination, VTOL, Space robotics.

\textbf{Почему:} Одно из ЛУЧШИХ совпадений вне топ-5. Она буквально занимается safe + resilient multi-agent UAV с математическими гарантиями. Лаба оснащена 30+ физическими дронами (Crazyflie, DJI, Hummingbird и др.). NSF CAREER, AFOSR YIP, NASA Early Faculty. Site Director NSF CAAMS.

\textbf{Совет:} Прочитай её статью про CBF + nonholonomic robots (2025). Напиши про safe multi-agent UAV с provably-correct гарантиями. Michigan = отличный «запасной план №1».
\end{gembox}

% ─── UPenn ───────────────────────────────────────────────────────────────────
\subsection{University of Pennsylvania}

\profsec{Vijay Kumar}{GRASP Lab}{88}
\textbf{Темы:} UAV swarms, Agile flight, Multi-robot coordination, Trajectory planning. TED Talk о квадрокоптерах.

\textbf{Почему:} Один из самых известных людей в мире по дронам. GRASP — один из старейших robotic labs в США.

\textbf{Совет:} GRASP очень конкурентный — нужен сильный проект с дронами и желательно один реальный эксперимент.

\profsec{Pratik Chaudhari}{GRASP / ESE}{68}
\textbf{Темы:} Deep learning theory, RL optimization, Physics-informed ML. Молодой профессор (2020).

% ─── UIUC ────────────────────────────────────────────────────────────────────
\subsection{UIUC — University of Illinois Urbana-Champaign}

\begin{starbox}{Naira Hovakimyan — ACRL (Advanced Controls Research Lab) \hfill \match{87}}
\textbf{Темы:} L1 Adaptive Control, Safe learning, UAV systems, NASA ULI robust autonomy, Nonlinear control.

\textbf{Почему:} Создательница L1 Adaptive Control — одного из мощнейших адаптивных регуляторов. Директор NASA University Leadership Initiative. UAV-работы тестировались на реальных самолётах и F-16. Если тебя цепляет математика управления — это твой человек.

\begin{tipbox}
\textbf{Совет:} Прочитай про L1 adaptive control перед контактом. Упомяни интерес к safe learning + неопределённость.
\end{tipbox}
\end{starbox}

% ─── UT Austin ───────────────────────────────────────────────────────────────
\subsection{UT Austin}

\profsec{Peter Stone}{LARG — Learning Agents Research Group}{75}
\textbf{Темы:} RL, Multi-agent systems, Cyber-physical systems. 30+ лет в RL. 2025 ACM/AAAI Allen Newell Award.

\profsec{Joydeep Biswas}{AMRL — Autonomous Mobile Robotics Lab}{65}
\textbf{Темы:} Autonomous navigation, Robot learning, Long-term autonomy.

% ─── UW ──────────────────────────────────────────────────────────────────────
\subsection{University of Washington}

\profsec{Byron Boots}{Robot Learning Lab}{78}
\textbf{Темы:} Robot learning, RL for dynamics, Sim-to-real, Off-road autonomy. Amazon Professor of ML. Бывший руководитель Guanya Shi.

% ─── CU Boulder ──────────────────────────────────────────────────────────────
\subsection{University of Colorado Boulder}

\profsec{Eric Frew}{RECUV — Research \& Engineering Center for Unmanned Vehicles}{80}
\textbf{Темы:} UAV systems, Autonomous flight, Planning under uncertainty, NSF CAAMS member.

\textbf{Почему:} Директор RECUV — одного из старейших UAV-centers в США. Работает с planning под uncertainty для дронов. Менее конкурентный при поступлении.

% ════════════════════════════════════════════════════════════════════════════
\section{УРОВЕНЬ 3 — Скрытые жемчужины и широкая сеть}
% ════════════════════════════════════════════════════════════════════════════

\subsection{UC Berkeley (дополнительные лабы)}

\begin{gembox}{Giuseppe Loianno — ARPL (Agile Robotics \& Perception Lab) \hfill \match{91}}
\textbf{Темы:} Agile drones, RL sim-to-real, Quadrotor control, Onboard learning, Swarms, RLtools library.

\textbf{Почему:} \textbf{ГОРЯЧИЙ МОМЕНТ — только что (2025) переехал из NYU в Berkeley!} Активно строит новую группу = много открытых мест. Лаборатория ARPL — одна из самых практических по RL для дронов. Соавтор RLtools (первый RL для дронов на микроконтроллере). NSF CAREER + DARPA YFA.

\begin{tipbox}
\textbf{Совет:} Написать СЕЙЧАС (через 1–2 года) — он только переехал. Показать код с gym-pybullet-drones на GitHub.
\end{tipbox}
\end{gembox}

\profsec{Koushil Sreenath}{Hybrid Robotics Lab}{89}
\textbf{Темы:} Aerial manipulation, CBF + RL для агрессивных манёвров дронов, Agile flight, Dynamic robots. Associate Professor. 16k+ цитирований.

\textbf{Совет:} Изучи его работы по CBF для aerial systems. Упомяни hybrid dynamics и safe flight.

\profsec{Mark Mueller}{High Performance Robotics Lab}{78}
\textbf{Темы:} Quadrotor control, Wind-resilient flight, Multi-UAV payload transport, Trajectory optimization.

\textbf{Почему:} Специализируется именно на квадрокоптерах. Меньше RL, больше математика управления. Меньшая лаба = больше внимания студентам.

\subsection{UCSD — UC San Diego}

\begin{gembox}{Sylvia Herbert — Safe Robotics Lab @ UCSD \hfill \match{85}}
\textbf{Темы:} Hamilton-Jacobi reachability analysis, Safe RL, Uncertainty quantification, Autonomous vehicles.

\textbf{Почему:} Один из ведущих экспертов по HJ-reachability — математическому фундаменту safe RL. Работает с Claire Tomlin (Berkeley) и Aaron Ames (Caltech). Молодой профессор (2021), активно берёт студентов. Тематика близка к Kochenderfer: safety + uncertainty.

\textbf{Совет:} Прочитай её статьи по HJ reachability. Напиши про formal safety guarantees + RL.
\end{gembox}

\profsec{Gaurav Sukhatme}{RESL / CAAMS}{70}
\textbf{Темы:} Networked robots, Field robotics, Marine UAV, Multi-agent. Директор CAAMS.

\subsection{Texas A\&M University}

\profsec{John Valasek}{Flight Dynamics \& Control Lab}{77}
\textbf{Темы:} Flight control, UAV systems, Adaptive control, Morphing aircraft, Real flight testing. Часть NSF CAAMS.

\profsec{Swaminathan Gopalswamy}{Unmanned Systems Lab}{65}
\textbf{Темы:} UAV systems, Autonomy, VTOL, Field testing. Директор Unmanned Systems Lab, сотрудничество с ВВС США.

\subsection{Virginia Tech}

\profsec{Mazen Farhood}{Advanced Control Systems Lab}{74}
\textbf{Темы:} Robust control, UAV verification, Formal methods, Safety analysis, Certification. Тематика верификации нейросетей совпадает с Kochenderfer.

\profsec{Craig Woolsey}{Nonlinear Systems Lab}{68}
\textbf{Темы:} Nonlinear control, UAV dynamics, Marine/aerial vehicles, Geometric mechanics.

\subsection{University of Minnesota}

\profsec{Nikolaos Papanikolopoulos}{LASR — Lab for Autonomous Systems}{70}
\textbf{Темы:} Small UAV, Computer vision, Autonomous navigation, Search and rescue. Работает с дронами с 1990-х.

\begin{gembox}{Peter Seiler — Control Science \& Dynamical Systems \hfill \match{72}}
\textbf{Темы:} Robust control, Safety certification, Aerospace, LPV systems. Работал с NASA и Boeing.

\textbf{Почему:} Если интересна safety certification для autonomous flight — хороший нишевый выбор.
\end{gembox}

\subsection{BYU — Brigham Young University}

\begin{gembox}{Randal Beard — MAGICC Lab \hfill \match{82}}
\textbf{Темы:} UAV formation control, Multi-agent cooperative systems, Planning, NSF CAAMS. 40\,000+ цитирований.

\textbf{Почему:} \textbf{Главная скрытая жемчужина.} BYU — не топ-10, но Beard — один из самых цитируемых людей в мире по UAV cooperative control. Автор классических учебников «Small Unmanned Aircraft» (с McLain). Часть NSF CAAMS (с Panagou, Michigan). Проще поступить, сильный руководитель.

\textbf{Совет:} Прочитай книгу «Small Unmanned Aircraft». Написать про cooperative UAV + safety planning.
\end{gembox}

\begin{gembox}{Timothy McLain — MAGICC Lab \hfill \match{78}}
\textbf{Темы:} UAV path planning, Cooperative flight, Real-time control, Hardware.

\textbf{Почему:} Соавтор Beard. Работают вместе в MAGICC Lab. Ко-адвайзор с Beard — можно написать к обоим.
\end{gembox}

\subsection{Ohio State University}

\profsec{Ran Dai}{Autonomous Control \& Optimization Lab}{72}
\textbf{Темы:} Autonomous aerospace, Trajectory optimization, Safety constraints, Neural control, UAV.

\subsection{Arizona State University}

\profsec{Heni Ben Amor}{Interactive Robotics Lab}{66}
\textbf{Темы:} Robot learning, RL from demonstrations, Transfer learning. Новая программа PhD in Robotics @ ASU (2024). Менее конкурентный входной барьер.

% ════════════════════════════════════════════════════════════════════════════
\section{Сводная таблица всех профессоров}
% ════════════════════════════════════════════════════════════════════════════

\begin{landscape}
\small
\begin{longtable}{p{3.2cm}p{3.5cm}p{2.2cm}p{5.5cm}p{1.2cm}p{1cm}p{1.5cm}}
\toprule
\textbf{Университет} & \textbf{Лаборатория} & \textbf{Профессор} &
\textbf{Ключевые темы} & \textbf{Совп.} & \textbf{★} & \textbf{💎} \\
\midrule
\endfirsthead
\toprule
\textbf{Университет} & \textbf{Лаборатория} & \textbf{Профессор} &
\textbf{Ключевые темы} & \textbf{Совп.} & \textbf{★} & \textbf{💎} \\
\midrule
\endhead
\bottomrule
\endfoot

\rowcolor{lightblue}
Stanford & SISL & Kochenderfer & MDP/POMDP, Safe RL, UAV, Julia & \match{98} & ★ & \\
Stanford & ASL & Pavone & Autonomous aerospace, Safety & \match{82} & & \\
\rowcolor{rowalt}
MIT & ACL & How & UAV, Multi-agent, POMDP & \match{90} & ★ & \\
\rowcolor{rowalt}
MIT & LIDS & Karaman & UAV, Algorithms, RRT* & \match{78} & & \\
\rowcolor{rowalt}
MIT & LIDS/MechE & Azizan & Adaptive control, Drones & \match{72} & & \\
CMU & LeCAR & Guanya Shi & RL+control, Aerial robotics & \match{88} & ★ & \\
CMU & AirLab & Scherer & UAV, Field robotics & \match{75} & & \\
\rowcolor{rowalt}
Caltech & CAST & Aaron Ames & CBF, Safety-critical & \match{80} & & \\
\rowcolor{rowalt}
Caltech & ARCL & Soon-Jo Chung & Aerospace robots, Multi-agent & \match{77} & & \\
UC Berkeley & HiPeR & Claire Tomlin & Reachability, Safe RL & \match{76} & & \\
UC Berkeley & RAIL & Sergey Levine & Deep RL, Robot learning & \match{65} & & \\
\rowcolor{rowalt}
Georgia Tech & Safe Robotics & Shreyas Kousik & Safe RL, Reachability & \match{85} & & 💎 \\
\rowcolor{rowalt}
Georgia Tech & LIDAR & Ye Zhao & RL + traj. opt., MPC & \match{72} & & \\
Michigan & DASC Lab & Dimitra Panagou & Safe UAV swarms, CBF & \match{92} & ★ & 💎 \\
UPenn & GRASP & Vijay Kumar & UAV swarms, Agile flight & \match{88} & ★ & \\
\rowcolor{rowalt}
UPenn & GRASP/ESE & Pratik Chaudhari & RL optimization & \match{68} & & \\
UIUC & ACRL & Naira Hovakimyan & L1 control, Safe learning & \match{87} & ★ & \\
UT Austin & LARG & Peter Stone & RL, Multi-agent & \match{75} & & \\
\rowcolor{rowalt}
UT Austin & AMRL & Joydeep Biswas & Autonomous navigation & \match{65} & & \\
UW & Robot Learning & Byron Boots & Robot learning, RL & \match{78} & & \\
\rowcolor{rowalt}
CU Boulder & RECUV/CAAMS & Eric Frew & UAV systems, Planning & \match{80} & & \\
UC Berkeley & Hybrid Robotics & Koushil Sreenath & CBF+RL, Aerial robotics & \match{89} & & 💎 \\
\rowcolor{lightgem}
UC Berkeley & ARPL & Giuseppe Loianno & Agile drones, RL sim-to-real & \match{91} & ★ & 💎 \\
UC Berkeley & HPR Lab & Mark Mueller & Quadrotor control, Wind & \match{78} & & \\
\rowcolor{rowalt}
UCSD & Safe Robotics UCSD & Sylvia Herbert & HJ-reachability, Safe RL & \match{85} & ★ & 💎 \\
\rowcolor{rowalt}
UCSD & RESL/CAAMS & Gaurav Sukhatme & Field robotics, UAV & \match{70} & & \\
Texas A\&M & FDCL & John Valasek & Flight control, Adaptive UAV & \match{77} & & \\
Virginia Tech & ACSL & Mazen Farhood & UAV verification, Robust & \match{74} & & \\
\rowcolor{rowalt}
Virginia Tech & Nonlinear Sys. & Craig Woolsey & Nonlinear dynamics, UAV & \match{68} & & \\
Minnesota & LASR & Papanikolopoulos & Small UAV, Navigation & \match{70} & & \\
\rowcolor{rowalt}
Minnesota & Control Sci. & Peter Seiler & Safety certification & \match{72} & & 💎 \\
\rowcolor{lightgem}
BYU & MAGICC Lab & Randal Beard & UAV formation, Cooperative & \match{82} & ★ & 💎 \\
\rowcolor{lightgem}
BYU & MAGICC Lab & Timothy McLain & Path planning, Real-time & \match{78} & & 💎 \\
Ohio State & ACOL & Ran Dai & Aerospace, Traj. opt. & \match{72} & & \\
\rowcolor{rowalt}
ASU & IRL & Heni Ben Amor & Robot learning, RL & \match{66} & & \\
\end{longtable}
\end{landscape}

% ════════════════════════════════════════════════════════════════════════════
\section{Стратегия поступления}
% ════════════════════════════════════════════════════════════════════════════

\subsection{Общий принцип}

\begin{stratbox}{Три правила PhD-поиска}
\begin{enumerate}[leftmargin=1.5em, itemsep=3pt]
  \item \textbf{Профессор важнее рейтинга университета.} Сильный руководитель в программе \#20 лучше слабого в \#5. Выбирай место где есть твой человек и тематика.
  \item \textbf{Подавать в 10–14 мест одновременно.} Это стандартная практика. Никто не подаёт только в одно.
  \item \textbf{PhD финансируется.} В США PhD в STEM практически всегда идёт с финансированием (stipend \$30–40k/год + tuition waiver). Это не аппликация «за свои деньги».
\end{enumerate}
\end{stratbox}

\subsection{4-летний roadmap}

\begin{center}
\begin{tabular}{>{\bfseries}p{2.8cm} p{3.5cm} p{8.5cm}}
\toprule
Период & Фокус & Конкретные задачи \\
\midrule
Годы 1--2 & Математика и теория &
  Линейная алгебра, ОДУ, теория вероятностей, MDP/POMDP-формализм, PID, LQR, устойчивость по Ляпунову. Читать: Kochenderfer «Algorithms for Decision Making», Sutton \& Barto RL. \\
\midrule
Годы 1--2 & Первый проект &
  PID vs RL vs Hybrid для квадрокоптера (gym-pybullet-drones + Stable-Baselines3). Добавить safety constraints (CBF) или POMDP-постановку. Выложить на GitHub. \\
\midrule
Год 3 & SISL-фактор &
  Добавить: POMDP (зашумлённые датчики), formal safety (CBF), или multi-agent coordination. Перевести часть кода на Julia + POMDPs.jl. \\
\midrule
Год 3--4 & Публикация &
  Workshop paper или full paper для ICRA, IROS, L4DC, ACC. Даже отказ — это опыт. Минимум — arXiv-препринт. Это главный дифференциатор. \\
\midrule
Год 4 & Нетворкинг &
  Посетить конференцию (ICRA, NeurIPS, CoRL). Найти профессоров из списка, представиться. Email: «Я прочитала статью X, вот что я сделала похожего [GitHub].» \\
\midrule
Год 4 (осень) & Подача документов &
  Декабрь — дедлайны большинства программ. Statement of Purpose: конкретный проект + конкретный профессор + конкретная исследовательская цель. \\
\bottomrule
\end{tabular}
\end{center}

\subsection{Рекомендуемый список для подачи (12 университетов)}

\begin{center}
\begin{tabular}{>{\bfseries}p{2.5cm} p{5.5cm} p{7cm}}
\toprule
Категория & Университет / Профессор & Обоснование \\
\midrule
\rowcolor{lightblue}
Мечта (3) & Stanford — Kochenderfer / Pavone &
  Главная цель. SISL = прямое совпадение. Сложно, но реально с сильным проектом. \\
\rowcolor{lightblue}
Мечта (3) & MIT — Jonathan How &
  UAV + multi-agent collision avoidance. Практически идеальное совпадение тем. \\
\rowcolor{lightblue}
Мечта (3) & CMU — Guanya Shi &
  RL + control + aerial robotics. Молодой проф, активно берёт студентов. \\
\rowcolor{lightgreen}
Сильные (4) & Michigan — Dimitra Panagou &
  Safe multi-agent UAV. 30+ дронов в лабе. Один из лучших outside топ-5. \\
\rowcolor{lightgreen}
Сильные (4) & UC Berkeley — Loianno / Sreenath &
  Loianno только переехал (2025) — много мест. Sreenath: CBF + agile flight. \\
\rowcolor{lightgreen}
Сильные (4) & Caltech — Aaron Ames &
  Если хочешь самую глубокую safety теорию — создатель CBF. \\
\rowcolor{lightgreen}
Сильные (4) & UIUC — Naira Hovakimyan &
  L1 adaptive control + NASA ULI. Лучший выбор при интересе к математике управления. \\
\rowcolor{lightyellow}
Реальные (3) & Georgia Tech — Shreyas Kousik &
  Safe RL + reachability. Молодой проф, меньше конкуренции чем в топ-5. \\
\rowcolor{lightyellow}
Реальные (3) & UCSD — Sylvia Herbert &
  HJ-reachability + safe RL. Молодой проф (2021), активно берёт студентов. \\
\rowcolor{lightyellow}
Реальные (3) & UPenn — Vijay Kumar &
  Легендарная UAV лаба. Нужен реальный эксперимент. \\
\rowcolor{lightorange}
Страховки (2) & BYU — Randal Beard &
  40k+ цит. по UAV cooperative control. Проще поступить, сильный руководитель. \\
\rowcolor{lightorange}
Страховки (2) & CU Boulder — Eric Frew &
  Специализированный UAV центр. CAAMS. Меньше конкуренция. \\
\bottomrule
\end{tabular}
\end{center}

\subsection{Как написать email профессору}

\begin{stratbox}{Структура идеального email}
\begin{description}[leftmargin=2em, style=nextline]
  \item[\textbf{Тема письма:}] \texttt{PhD inquiry — [Название статьи которую ты прочитала]}\\
    Конкретная тема сразу говорит: «Я читала твою работу, это не массовая рассылка.»
  \item[\textbf{Абзац 1:}] Кратко кто ты и откуда. 1–2 предложения. Университет, специальность, год.
  \item[\textbf{Абзац 2:}] Конкретная статья + вопрос или наблюдение. Покажи что реально читала.\\
    \textit{Пример:} «В вашей статье RAIL вы показываете, что reachability-based safety layer можно добавить к любой RL-политике. Я задалась вопросом: применимо ли это к POMDP-постановке где наблюдения зашумлены?»
  \item[\textbf{Абзац 3:}] Твой проект с ссылкой на GitHub. Конкретный результат.\\
    \textit{Пример:} «Я реализовала сравнение PID vs PPO для квадрокоптера с ветровыми возмущениями. Репозиторий: [ссылка]. Планирую добавить CBF safety layer.»
  \item[\textbf{Финал:}] «Планируете ли вы брать PhD-студентов в [осень 20XX]?»\\
    Прямой вопрос уважает время профессора.
\end{description}
\end{stratbox}

\vspace{4pt}

\begin{tipbox}
\textbf{Золотое правило:} Профессора получают десятки «cold emails» в неделю. Тебя выделит: конкретная статья + конкретный вопрос + ссылка на реальный код. Без этих трёх элементов письмо теряется.
\end{tipbox}

\subsection{Statement of Purpose: структура}

\begin{enumerate}[leftmargin=1.5em, itemsep=4pt]
  \item \textbf{Хук (1 абзац):} Конкретная исследовательская проблема которую ты хочешь решить. Не «я мечтаю о роботах», а «Collision avoidance для autonomos UAV в условиях неопределённости — нерешённая проблема которую я хочу формализовать как POMDP и решить методами safe RL.»
  \item \textbf{Твой фон (1--2 абзаца):} Конкретные проекты, курсы, математический бэкграунд. Цифры и результаты важнее названий.
  \item \textbf{Твой проект (1 абзац):} PID vs RL vs Hybrid + что ты добавила (POMDP? CBF?). Ссылка на GitHub или arXiv.
  \item \textbf{Почему этот профессор (1 абзац):} Конкретная статья + почему именно она. Не «у вас сильная лаба», а «ваш подход к CBF в работе X именно та математика которую я хочу использовать для задачи Y.»
  \item \textbf{Исследовательский вопрос (1 абзац):} Что конкретно ты хочешь исследовать в PhD. Не обязательно точно — но конкретно в направлении.
\end{enumerate}

\vspace{6pt}
\begin{center}
{\small\color{darkgray}\itshape
Документ составлен на основе исследовательского трека: PhD @ Stanford SISL.
Все данные актуальны на апрель 2026.}
\end{center}

\end{document}